% =====================================================================
% Exemple minimal — DIAPOSITIVES
% Le support est déduit de la classe : beamer suffit, aucune option à passer.
% =====================================================================
\documentclass[9pt,t]{beamer}
\usepackage[lang=fr, theme=n7-dark, math=analysis,
            institution=n7,insa]{ocots}

\title{Exemple de diapositives}
\author{Prénom \textsc{Nom}}
\date{\today}

\begin{document}

\slidechapter{1}{Titre du chapitre}
\slidetitlepage% remplace les 28 lignes recopiées dans chaque jeu de la v0

% ---------------------------------------------------------------- contenu
\begin{slide}{Titre de la diapositive}

Texte, avec une norme $\norm{x}$.

\begin{definition}{Une définition}{}
    Énoncé.
\end{definition}

\pause

\begin{remark}
    Une remarque.
\end{remark}

\end{slide}

% La couleur d'en-tête peut changer par thème abordé.
\begin{slide}[\ocotscolor{slide2}]{Une autre couleur d'en-tête}

\begin{proofbegin}
    Début d'une preuve étalée sur plusieurs diapositives.
\end{proofbegin}

\end{slide}

\begin{slide}{Suite de la preuve}

\begin{proofend}
    Fin de la preuve.
\end{proofend}

% <<< étape 3 : tous les environnements, toutes les options >>>

\end{slide}

\end{document}
